\section{Related Work}
\subsection{Compact Vision Models}
Residual networks \cite{he2016}, VGG \cite{simonyan2015}, EfficientNet \cite{tan2019}, and MobileNetV2 \cite{sandler2018} represent different routes to high accuracy or reduced computation. The parameter count alone, however, does not determine batch-one CPU latency: depthwise operations, feature-map size, and implementation details also affect execution \cite{kyrkou2026}. Although lightweight vision transformers have narrowed the gap with CNNs on general benchmarks \cite{mehta2022}, they tend to underperform convolutional architectures in data-constrained, sub-1M-parameter settings \cite{kyrkou2026}. For medical images, fine-tuned pretrained backbones provide strong reference points \cite{tajbakhsh2016}, but a comparison restricted to similarly large pretrained models cannot establish what is lost by imposing a much tighter memory budget. We therefore contrast pretrained, fine-tuned backbones with scratch-trained compact counterparts rather than treating their training regimes as equivalent.

Kyrkou's LS-Net was designed for the sub-1M vision regime using multi-path atrous convolutions \cite{kyrkou2026}. Parallel depthwise paths with distinct dilation factors enlarge the set of spatial contexts represented without requiring a proportionate increase in learned parameters. In the present work, LS-Net serves as an adapted architectural baseline; our contribution is the inference strategy, not the backbone. The new question is whether such an architecture remains useful across two MRI tasks and whether its errors can be selectively addressed at inference time. A compact MPAC-ResNet variant helps separate the value of parallel atrous processing from the broader architectural choices of LS-Net.

\subsection{Distribution Shift and Test-Time Augmentation}
A high score on a balanced collection of two-dimensional images cannot establish robustness to different MR sequences, clinical labels, or subject composition \cite{guan2022}. In BraTS, HGG and LGG are grades rather than the tumor-type labels of the first benchmark \cite{menze2015}. We therefore report task-specific results and do not present the two test sets as interchangeable estimates of one distribution. Splitting by subject avoids putting different slices from the same patient in both training and test sets, a leakage mechanism documented in brain MRI classification \cite{yagis2021}. Patient-level aggregation also prevents repeated slices of one subject from being mistaken for independent clinical cases.

Test-time augmentation averages predictions from transformed versions of the same input, an approach studied in medical imaging as well as general image classification \cite{wang2019,moshkov2020,shanmugam2021}. Recent generative approaches use diffusion models to produce domain-aware augmented views at test time \cite{ma2026}, but their computational overhead is substantially higher than deterministic transforms. Unconditional use improves coverage at a compute cost on every slice, but can also overturn correct predictions \cite{shanmugam2021}; confidence gating spends compute on a smaller subset, but confidence alone cannot identify which transform is informative. A softmax score is also not necessarily a calibrated estimate of correctness \cite{guo2017}. Our method couples a low-confidence gate to a pair of scale probes that are motivated by LS-Net's multi-scale feature paths. Disagreement with the original prediction determines whether to acquire the remaining views. This is an inference-time intervention, not a change to training, calibration, or network weights.
